\chapter{Limity a Derivace} \label{ch:Limity a Derivace}
\section{O co by se mělo jednat}
\section{O co se jedná}
\begin{definice}[Funkce] \label{def:Funkce} Funkce (= zobrazení) se skládá z následujících dat:
\begin{itemize}
	\item Množiny $A$, zvané \emph{definiční obor},
	\item Množiny $B$, zvané \emph{obor hodnot},
	\item Pravidla, které každému prvku $a\in A$ přiřadí právě jeden prvek $f(a)\in B$. Tomuto pravidlu říkáme předpis funkce.
\end{itemize}
Značení $f\colon A\to B$ znamená, že $A$ je definiční obor funkce $f$ a $B$ je její obor hodnot.
\end{definice}
\newcommand{\Leden}{\textnormal{Leden}}
\newcommand{\Unor}{\textnormal{Únor}}
\newcommand{\Brezen}{\textnormal{Březen}}
\newcommand{\Prosinec}{\textnormal{Prosinec}}
\begin{příklad}
	$A=\{\Leden, \Unor, \Brezen, ..., \Prosinec\}$, $B=\{28,30,31\},$
	\begin{align*}
		f(\Leden)&=31 \\
		f(\Unor)&=28 \\
		f(\Brezen)&=31 \\
			 &\vdots
	\end{align*}
Funkce $f\colon A\to B$ udává počet dnů měsíce v nepřestupném roce. Není daná žádným matematickým vzorečkem, její definiční obor neobsahuje čísla. 
\end{příklad}

\begin{poznámka}
\leavevmode
\begin{itemize}
	\item Definiční obor a obor hodnot jsou \emph{součástí definice} a mají vliv na vlastnosti funkce. Například funkce $f\colon (0,\infty)\to \mathbb{R}, \ f(x)=x^{2}$ je rostoucí, zatímco $g\colon \mathbb{R}\to \mathbb{R}, \ g(x)=x^{2}$ nikoliv.
	\item Může se stát, že dva různé prvky $a_1\neq a_2\in A$ funkce přiřadí na stejný prvek $f(a_1)=f(a_2).$ Také se může stát, že existuje prvek $b\in B$, který není obrazem žádného $a\in A$, tedy
	\begin{equation*}
	  \forall a\in A\colon b\neq f(a).
	\end{equation*}
	Například funkce $g$ z předchozího bodu má obě tyto vlastnosti:
	\begin{align*}
		3^{2} &= (-3)^{2} \\
		\forall a\in \mathbb{R}\colon-1&\neq a^{2}
	\end{align*}
\end{itemize}
\end{poznámka}
Dále nebudeme značit definiční obory a obory hodnot, pokud to nebude nutné. Zpravidla budeme uvažovat největší možné definiční obory v rámci $\mathbb{R}.$
\begin{definice}[Limita funkce] \label{def:Limita funkce}
	Nechť $c\in \mathbb{R}$ a $f$ je funkce, jejíž obor hodnot je $\mathbb{R}$ a její definiční obor obsahuje $(c-\delta,c+\delta)\setminus\{c\}$ pro nějaké $\delta>0$. Limita funkce $f$ v bodě $c$, pokud existuje, je číslo $L$ splňující
\begin{equation*}
	\forall \epsilon>0 \ \exists \delta>0\ \forall x\in (c-\delta,c+\delta)\setminus\{c\}\colon f(x)\in(L-\epsilon, L+\epsilon).
\end{equation*}
Značíme ho
\begin{equation*}
	\lim_{x\to c}f(x):=L,
\end{equation*}
\end{definice}

\begin{úloha}
	Nechť $f(x)=2x+3.\ \lim_{x\to 1}f(x)=\ ?$
\end{úloha}

\begin{úloha}
  Nechť $f\colon \mathbb{R}\to \mathbb{R}$ má předpis
  \begin{equation*}
	  f(x)=\left\{
	\begin{matrix}
		2, & x\neq 0 \\
		0  & x=0
	\end{matrix}\right.
  \end{equation*}
  A $g\colon \mathbb{R}\setminus \{0\}\to \mathbb{R}$ má předpis $g(x)=2.$ $\lim_{x\to 0}f(x),\ \lim_{x\to 0}g(x)=\ ?$
\end{úloha}

\begin{příklad}
	Nechť $f(x)=x^{2}. \lim_{x\to 3}f(x)=\ ?$
\end{příklad}

\begin{příklad}
  Nechť $f(x)=\left\{
  \begin{matrix}
	  0, & x\leq 0 \\
	  1, & x>0
  \end{matrix}\right.$\par $\lim_{x\to 0}f(x)=\ ?$
\end{příklad}

Vidíme, že pokud se nesnažíme konstruovat funkce s patologickým předpisem, nebo násilím odebranými body z definičního oboru, většinou platí $\lim_{x\to c}f(x)=f(c).$

\begin{definice}[Spojitá funkce]
  Mějme funkci $f\colon (a,b)\to \mathbb{R},\ x\in (a,b).$ Říkáme, že $f$ je spojitá v bodě $x\in (a,b)$, pokud platí
  \begin{equation*}
	  \lim_{y\to x}f(y)=f(x), \text{ neboli } \lim_{h\to 0}f(x+h)=f(x).
  \end{equation*}
  Říkáme, že $f$ je spojitá na $(a,b),$ pokud je spojitá v bodě $x$ pro všechna $x\in (a,b).$
\end{definice}

\begin{Todo}
  polynomy jsou spojite
\end{Todo}

\begin{definice}[Derivace funkce]
  Mějme funkci $f\colon (a,b)\to \mathbb{R},\ x\in (a,b).$ Derivace funkce $f$ v bodě $x$ (pokud existuje) je číslo
  \begin{equation*}
	  f'(x) := \lim_{h\to 0}\frac{f(x+h)-f(x)}{h}
  \end{equation*}
  Pokud existuje $f'(x)$ pro všechna $x\in (a,b)$, značíme odpovídající funkci $f'$ nebo také $\tfrac{df}{dx}:= f'.$
\end{definice}

\begin{úloha}
  $f(x)=C$ (konstanta). $f'(x)=\ ?$
\end{úloha}
\begin{úloha}
  $f(x)=ax+b.\ f'(x)=\ ?$ 
\end{úloha}
\begin{úloha}
  $f(x)=x^{2}.\ f'(x)=\ ?$ 
\end{úloha}
\begin{příklad}
  $f(x)=x^{n}, n\geq 2.\ f'(x)=\ ?$\par
\textit{Řešení:} Roznásobením závorek v $(x+h)^{n}$ dostáváme členy typu $x^{a}h^{b},$ $a+b=n.$
\begin{equation*}
	(x+h)^{n} = (x+h)\cdot \hdots \cdot (x+h) = x^{n}+nx^{n-1}h + (\hdots)h^{2}+\hdots + h^{n}
\end{equation*}
Členy $x^{a}h^{b}, b\geq 2$ nás příliš netrápí, protože k limitě $\lim_{h\to 0}\frac{(x+h)^{n}-x^{n}}{h}$ nepřispějí. Důležité je, že u členu $x^{n-1}h$ je koeficient $n$, protože tolik existuje možností, jak si z jedné závorky \uv{vybrat} $h$ a z ostatních $x$. Tedy
\begin{equation*}
	(x^{n})'=\lim_{h\to 0}\frac{(x+h)^{n}-x^{n}}{h}=\lim_{h\to 0}\frac{nx^{n-1}h+(\hdots)h^{2}+\hdots+h^{n}}{h}=nx^{n-1}
\end{equation*}
\end{příklad}

\begin{definice}[Součet, násobek a součin funkcí]
  Nechť $\alpha\in \mathbb{R},\ f\colon A\to \mathbb{R},\ g\colon A\to \mathbb{R}.$ Pak definujeme funkce $f+g, f\cdot g, \alpha f$ tak, že
  \begin{itemize}
	  \item $f+g\colon A\to \mathbb{R},\ f\cdot g\colon A\to \mathbb{R},\ \alpha f\colon A\to \mathbb{R}$
	  \item $(f+g)(x):= f(x)+g(x),\ (f\cdot g)(x):= f(x)\cdot g(x),\ (\alpha f)(x):= \alpha\cdot f(x)$
  \end{itemize}
\end{definice}

\begin{Todo}
  motivace leibnize
\end{Todo}

\begin{věta}[Derivace součtu, součinu a násobku]
  Nechť $f,g\colon (a,b)\to \mathbb{R},\ x\in (a,b), \alpha\in \mathbb{R}$. Pak platí
  \begin{align}
	  (f+g)'(x) &=f'(x)+g'(x),\label{eq:derivace souctu} \\
	  (f\cdot g)'(x)&=f(x)g'(x)+f'(x)g(x), \label{eq:derivace soucinu} \\
	  (\alpha f)'(x)&=\alpha f'(x) \label{eq:derivace nasobku}
  \end{align}
  existují-li čísla napravo (tedy existují-li $f'(x),\ g'(x)$).
\end{věta}
\begin{důkaz} Rovnost \ref{eq:derivace souctu} se dokáže následovně:
  \begin{align*}
	  (f+g)'(x)&=\lim_{h\to 0}\frac{(f+g)(x+h)-(f+g)(x)}{h}\\
		   &=\lim_{h\to 0}\frac{f(x+h)-f(x)+g(x+h)-g(x)}{h}\\
		   &=\lim_{h\to 0}\frac{f(x+h)-f(x)}{h}+\lim_{h\to 0}\frac{g(x+h)-g(x)}{h}=f'(x)+g'(x).
  \end{align*}
  Rovnost \ref{eq:derivace soucinu} plyne z výpočtu
  \begin{align*}
	  (fg)'(x)&=\lim_{h\to 0}\frac{(fg)(x+h)-(fg)(x)}{h}\\
		  &=\lim_{h\to 0}\frac{f(x+h)g(x+h)-f(x)g(x)}{h}\\
		  &=\lim_{h\to 0}\frac{f(x+h)g(x+h)-f(x+h)g(x)+f(x+h)g(x)-f(x)g(x)}{h}\\
		  &=\lim_{h\to 0}f(x+h)\frac{g(x+h)-g(x)}{h}+\lim_{h\to 0}\frac{f(x+h)-f(x)}{h}g(x)\\
		  &=(\lim_{h\to 0}f(x+h))\left(\lim_{h\to 0}\frac{g(x+h)-g(x)}{h}\right)+\left(\lim_{h\to 0}\frac{f(x+h)-f(x)}{h}\right)g(x)\\
		  &= f(x)g'(x)+f'(x)g(x)
  \end{align*}
  Využili jsme toho, že pokud $f$ má v $x$ derivaci, je tam spojitá. To plyne z výpočtu
  \begin{align*}
	  \lim_{h\to 0}f(x+h) &= f(x)+\lim_{h\to 0}\left(\frac{f(x+h)-f(x)}{h}h\right) \\
			      &= f(x)+ \lim_{h\to 0}\left(\frac{f(x+h)-f(x)}{h}\right)\lim_{h\to 0}h \\
			      &= f(x) + f'(x)\cdot 0 = f(x)
  \end{align*}
  Rovnost \ref{eq:derivace nasobku} plyne z \ref{eq:derivace soucinu}, pokud za $g$ vezmeme konstantní funkci.
\end{důkaz}

\begin{příklad}
  Najděte rovnici tečny ke grafu funkce $f(x)=x^{2}-4x+3$ v bodě $(1,f(1))$.\par
\textit{Řešení:} Derivace $f'(x)=2x-4$ má v bodě $x=1$ hodnotu $f'(1)=-2,$ což bude směrnice tečny. Tečna je popsána rovnicí $y=-2x+b$ a prochází bodem $(x,y)=(1,f(1))=(1,0),$ z toho $b=2.$ Tečna je tedy popsána rovnicí $y=-2x+2.$
\end{příklad}
\begin{příklad}
  $f(x)=\sqrt{x}.\ f'(x)=\ ?$\par
\textit{Řešení:}
\begin{align*}
	\sqrt{x}\sqrt{x}   &= x \hspace{3em} \rvert \frac{d}{dx} \\
	\sqrt{x}(\sqrt{x})'+(\sqrt{x})'\sqrt{x} &= 1 \\
	2 \sqrt{x} (\sqrt{x})' &= 1 \\
	(\sqrt{x})' &= \frac{1}{2 \sqrt{x}} \\
	\textnormal{neboli } \left(x^{\frac{1}{2}}\right)' &= \frac{1}{2}x^{\frac{1}{2}-1}
\end{align*}
Vidíme, že ač za definiční obor pro $f(x)=\sqrt{x}$ můžeme vzít $[0,\infty)$, derivace existuje jen na $(0,\infty).$
\end{příklad}

\begin{úloha}
  $f(x)=\frac{1}{x}.\ f'(x)=\ ?$
\end{úloha}

\begin{definice}[Složená funkce]
  Nechť $f\colon A\to B$ a $g\colon B\to C$ jsou funkce. Pak složená funkce $g\circ f$ má
\begin{itemize}
\item definiční obor $A$ a obor hodnot $C$,
\item předpis $(g\circ f)(x) := g(f(x))$
\end{itemize}
\end{definice}
\begin{Todo}
  komutativni diagram
\end{Todo}

\begin{příklad}\label{pr:slozena funkce}
	$f\colon [-\frac{5}{9},\infty)\to [0,\infty),\ f(x)=9x+5,\ g\colon [0,\infty)\to [0,\infty),\ g(x)=\sqrt{x}.$ Složená funkce je $g\circ f\colon [-\frac{5}{9},\infty)\to [0,\infty)$ s předpisem $(g\circ f)(x)=\sqrt{9x+5}$
\end{příklad}

\begin{věta}[Derivace složené funkce]
  Nechť $f\colon A\to B$ a $g\colon B\to C,$ kde $A,B,C\subset \mathbb{R}.$ Pak pro $x\in A$ platí
  \begin{equation*}
	  (g\circ f)'(x)=g'(f(x))\cdot f'(x),
  \end{equation*}
  existuje-li číslo napravo (tedy existují-li $g'(f(x))$ a $f'(x)$).
\end{věta}
\begin{náčrt důkazu}
  \begin{align*}
	  (g\circ f)'(x) &= \lim_{h\to 0}\frac{g(f(x+h))-g(f(x))}{f(x+h)-f(x)}\frac{f(x+h)-f(x)}{h} \\
			 &= \lim_{h\to 0}\left(\frac{g(f(x+h))-g(f(x))}{f(x+h)-f(x)}\right)\lim_{h\to 0}\left(\frac{f(x+h)-f(x)}{h}\right) \\
			 &= \lim_{k\to 0}\left(\frac{g(f(x)+k)-g(f(x))}{k}\right)\lim_{h\to 0}\left(\frac{f(x+h)-f(x)}{h}\right) \\
			 &= g'(f(x))f'(x)
  \end{align*}
  Na dokončení důkazu bychom museli vysvětlit, proč jde na prvním řádku rozšířit zlomek číslem $f(x+h)-f(x)$ (číslo musí být nenulové), a že můžeme od limity $h\to 0$ přejít substitucí $k:=f(x+h)-f(x)$ na limitu $k\to 0.$
\end{náčrt důkazu}

\begin{příklad}[Pokračování příkladu \ref{pr:slozena funkce}]
  $h(x)=\sqrt{9x+5}.\ h'(x)= ?$\par
  \textit{Řešení:} $h=g\circ f$ z příkladu \ref{pr:slozena funkce}, takže 
\begin{equation*}
  (\sqrt{9x+5})' = h'(x)=g'(f(x))\cdot f'(x)=\frac{1}{2 \sqrt{9x+5}}\cdot 9 = \frac{9}{2 \sqrt{9x+5}}
\end{equation*}
\end{příklad}

\begin{příklad}
	Najděte rovnici tečny ke kružnici $x^{2}+y^{2}=\textcolor{red}{25}$ v bodě $(\textcolor{blue}{3},\textcolor{violet}{4}).$\par
  \textit{Řešení}: stačí místo kružnice uvažovat graf funkce $h(x)=\sqrt{25-x^{2}}.$ Ta se napíše jako $h=g\circ f,\ g(x)=\sqrt{x}, f(x)=25-x^{2}.$ Z toho máme
  \begin{equation*}
	  (\sqrt{25-x^{2}})'=h'(x)=g'(f(x))\cdot f'(x)=\frac{1}{2 \sqrt{25-x^{2}}}\cdot (-2x) = \frac{-x}{\sqrt{25-x^{2}}}.
  \end{equation*}
  Protože $h'(3)=-\frac{3}{4},$ tečna má rovnici $y=-\frac{3}{4}x+b.$ Dosazením bodu $(3,4)$ dostaneme $b=\frac{25}{4}.$ Rovnice jde napsat v elegantním tvaru $\textcolor{blue}{3}x+\textcolor{violet}{4}y=\textcolor{red}{25}.$
\end{příklad}

\begin{věta*}[Vlastnosti funkce plynoucí z derivace]
  Nechť $f\colon (a,b)\to \mathbb{R}$ má na $(a,b)$ derivaci. Pak
  \begin{itemize}
	  \item Pokud $f'>0$ na $(a,b)$, je $f$ rostoucí na $(a,b)$
	  \item Pokud $f'<0$ na $(a,b)$, je $f$ klesající na $(a,b)$
	  \item Pokud má $f$ lokální extrém v bodě $x\in (a,b)$, platí $f'(x)=0.$
	  \item Pokud $f'(x)=0$ a existuje $\epsilon>0$ takové, že $f$ je klesající na $(x-\epsilon,x)$ a rostoucí na $(x,x+\epsilon)$, má $f$ v bodě $x$ lokální minimum. Podobně pro lokální maximum.
  \end{itemize}
\end{věta*}

\begin{úloha}
  Najděte souřadnice vrcholu paraboly, jejíž rovnice je $y=3x^{2}-6x+7$
\end{úloha}

\begin{úloha}
  Načrtněte graf funkce $f(x)=x^{3}-\frac{3}{2}x^{2}-6x+2$ tak, že najdete lokální extrémy a intervaly, kde je funkce rostoucí/klesající.
\end{úloha}
